\subsection{Residuals and Scaling Relations}
  \label{subsec:residuals}

  We quantify the quality of the fits using radial velocity residuals
  \begin{equation}
    \Delta v(r) = v_{\mathrm{obs}}(r) - v_{\mathrm{eff}}(r) .
  \end{equation}

  The residuals exhibit no systematic radial trends across the sample.
  In particular, no characteristic scale or radius-dependent bias is observed.

  The effective model naturally reproduces the empirical correlation between
  rotation curve shape and baryonic surface density.
  Galaxies with higher central surface density remain in the Newtonian regime
  over a larger radial range, while diffuse systems enter the saturated regime earlier.

  The transition radius at which rotation curves flatten corresponds closely to
  the regime where the effective surface flux density approaches the saturation limit.
  Beyond this point, further decreases in baryonic density no longer translate
  into stronger dynamical suppression, leading to asymptotically flat rotation velocities.

  This behavior is closely related to the observed radial acceleration relation.
  In the present framework, this relation is not imposed but emerges directly
  from the effective saturation of the gravitational response at low acceleration~\cite{McGaugh2016RAR}.
